\originalchapter{作业八：转动、吸附与恒星的经典障碍}
\psource{8.044，2013 春，PS8，4 月 10 日截止。中文题面；据官方答案核对的 AI 独立解答。}
\originalsection{1．空间针状尘粒}
\begin{exercise}
星际尘粒可呈针状。原题简化主轴转动模型用三组角与共轭动量 $\theta_i,L_i$，$H=\sum_{i=1}^3L_i^2/(2I_i)$，平衡温度 $T$。
(a) 给无未算积分的六变量联合密度；(b) 若轴 3 沿长轴且 $I_3\ll I_1=I_2$，动量更可能平行还是垂直长轴，定量说明；(c) $N$ 颗尘粒的转动熵；(d) 是否满足第三定律；(e) 是否出现能隙行为？
\end{exercise}
\begin{solution}
按原题把三个角作为各周期 $2\pi$ 的独立规范坐标，$p=(2\pi)^{-3}\prod_i(2\pi I_i\kb T)^{-1/2}e^{-L_i^2/(2I_i\kb T)}$。(b) $\avg{L_i^2}=I_i\kb T$，横向平方与纵向平方之比 $(I_1+I_2)/I_3\gg1$，故趋向垂直。(c) 单粒子 $z_R=(2\pi)^{9/2}(I_1I_2I_3)^{1/2}(\kb T)^{3/2}/h^3$；对以各尘粒位置区分的转动自由度，$S_R=N\kb(\ln z_R+3/2)$。(d) $T\to0$ 时趋负无穷，经典近似失效，须用量子谱。(e) 经典能量连续，无低温激活能隙；量子有限转子则有。一般真实刚体的 Euler 角不都是独立周期规范变量，需正确 $SO(3)$ 测度与动量；这里不把题设的简化角体积冒称一般刚体结果。
\end{solution}
\originalsection{2．台阶表面的吸附}
\begin{exercise}
硅表面 $M$ 个吸附位中，1\% 为拐角位、能量零，14\% 为边缘位、能量 $\Delta$，85\% 为面位、能量 $1.5\Delta$。$N$ 个相同氙原子、$M\gg N$，可忽略占位竞争及动能。
(a) 求 $Z(N,T)$；(b) 面位与角位数之比；(c) $\kb T\ll\Delta$ 热容；(d) 高温面位概率；(e) $T\to\infty$ 熵；(f) 是否有能隙行为？
\end{exercise}
\begin{solution}
令 $z=M[.01+.14e^{-\beta\Delta}+.85e^{-3\beta\Delta/2}]$。相同粒子的稀占位计数为 $Z\simeq z^N/N!$；若临时给原子编号，条件概率用 $z^N$ 也相同，但绝对熵不同。官方答案使用编号式，(e) 未扣置换；此处延续作业五的正确计数。
(b) 比值 $85e^{-3\beta\Delta/2}$。(c) $U\simeq14N\Delta e^{-\beta\Delta}$，$C\simeq14N\kb(\beta\Delta)^2e^{-\beta\Delta}$，需同时满足 $14e^{-\beta\Delta}\ll1$。(d) 高温概率 $.85$。(e) $S\simeq\kb[N\ln M-\ln N!]$；严格单占位模型高温为 $\kb\ln\binom MN$，稀占位时一致，编号式为 $N\kb\ln M$。(f) 最低激发间隔 $\Delta$，出现激活指数热容，但角位仍有组态简并。
\end{solution}
\originalsection{3．线性势阱的绝热冷却}
\begin{exercise}
$N$ 个不可分辨经典非相互作用原子在三维阱 $V(r)=ar$ 中，温度 $T$。
(a) 求 $z=AT^\alpha a^{-\eta}$ 的常数与指数；(b) 用 $N,\kb,z$ 表示无导数的熵；(c) 无热交换可逆降低 $a$，求 $T(a)$，初值 $T_0,a_0$。
\end{exercise}
\begin{solution}
动量积分 $(2\pi m\kb T)^{3/2}$，位置积分 $4\pi\int r^2e^{-\beta ar}\dd r=8\pi(\kb T/a)^3$，因此 $A=8\pi\kb^3(2\pi m\kb)^{3/2}/h^3$、$\alpha=9/2,\eta=3$。$Z=z^N/N!$，$U=9N\kb T/2$，$S=\kb[N\ln z-\ln N!+9N/2]\simeq N\kb[\ln(z/N)+11/2]$。等熵即 $z$ 固定，$T=T_0(a/a_0)^{2/3}$。
\end{solution}
\originalsection{4．二维氢分子气体}
\begin{exercise}
面积 $A$ 上 $N$ 个非相互作用分子，转动限于平面，量子数 $m\in\mathbb Z$，每个 $m$ 一态，$\epsilon_m=\hbar^2m^2/(2I)$。本题不另加核自旋对称性限制。
(a) 写出单分子转动配分和，不须解析化；(b) 求 $p(3)/p(2)$；(c) 求 $p(m=3\mid\epsilon=9\hbar^2/(2I))$ 和 $p(m=1\mid\epsilon\le\hbar^2/(2I))$；(d) 高温转动内能。
\end{exercise}
\begin{solution}
令 $\epsilon_1=\hbar^2/(2I)$，$z_R=\sum_{m=-\infty}^{\infty}e^{-\beta\epsilon_1m^2}=1+2\sum_{j\ge1}e^{-\beta\epsilon_1j^2}$。比值 $e^{-5\beta\epsilon_1}$。两个条件概率分别 $1/2$、$1/(2+e^{\beta\epsilon_1})$。高温和变积分 $z_R\simeq\sqrt{\pi/(\beta\epsilon_1)}$，$U_R=N\kb T/2$，一个二次转动自由度。
\end{solution}
\originalsection{5．恒星为何需要隧穿}
\begin{exercise}
按原题的经典检验，质子电荷 $e=4.8\times10^{-10}$ esu 均匀在半径 $R=1.2\times10^{-13}$ cm 球内，另一质子作点电荷。
(a) 进入距离 $R$ 所需最小能量；(b) 三维理想气体 $T=4.0\times10^7$ K 中能量超过它的概率 $p_+$，假定阈值远大于 $\kb T$，可用 $\int_a^\infty\sqrt y e^{-y}\dd y\simeq\sqrt a e^{-a}$；(c) 核碰撞截面 $\sigma=\pi(2R)^2$，中心质量密度 $100\,\mathrm{g\,cm^{-3}}$，求平均速率、平均自由程 $L=1/(n\sigma)$、自由时间；(d) 每次碰撞融合概率 $p_+$ 独立，求每质子等待时间，与题设宇宙年龄 150 亿年比较；太阳质量 $2\times10^{33}$ g，假定全有此中心条件，估计每秒融合事件数。历史背景为经典高能尾无法解释持续核聚变，后来需要穿透 Coulomb 势垒的量子机制；外书长引文未整段转载。这里的温度、宇宙年龄仅保留题设数值，不作为当前测量。
\end{exercise}
\begin{solution}
(a) 球外势如点电荷，$E_{\min}=e^2/R=1.92\times10^{-6}$ erg，约 1.20 MeV。
(b) $y=E_{\min}/(\kb T)\simeq347.6$，$p_+\simeq(2/\sqrt\pi)\sqrt y e^{-y}$。以原解常数 $\kb=1.381\times10^{-16}$ erg/K、$m_p=1.67\times10^{-24}$ g 计算，$\log_{10}p_+\simeq-149.627$，$p_+\simeq2.36\times10^{-150}$；原题正文预告 $10^{-148}$ 与自身数字不一致，须以实算为准。
(c) $\avg v=\sqrt{8\kb T/(\pi m_p)}\simeq9.18\times10^7$ cm/s，$\sigma\simeq1.810\times10^{-25}$ cm$^2$、$n\simeq5.988\times10^{25}$ cm$^{-3}$，$L\simeq.0923$ cm，$\tau\simeq1.01\times10^{-9}$ s。
(d) $\tau_f=\tau/p_+\sim4.3\times10^{140}$ s，远大于 $4.73\times10^{17}$ s；$N_\odot=2\times10^{33}/m_p$。每秒每质子参与次数 $Np_+/\tau$ 为约 $2.8\times10^{-84}$，而\textbf{不同双质子碰撞事件}须除二，约 $1.4\times10^{-84}$ s$^{-1}$。官方答案未区分参与次数与双体事件。经典越垒机制完全不足；实际 pp 链还需弱相互作用与空间结构，隧穿本身不能直接给总融合率。
\end{solution}
